\honest{Conditional Independence, and Naive Bayes}{Eliezer Yudkowsky}{2008}

Last time I defined mutual information: $I(X;Y) = H(X) + H(Y) - H(X,Y)$. My example had X with eight equally likely states and Y with four, so 3 and 2 bits of entropy, and the rule that X and Y are both even or both odd. That leaves 16 joint states, 4 bits, and 1 bit of mutual information.

Now add Z, the question ``Are X and Y even or odd?'' Z has 1 bit of entropy, and 1 bit of mutual information with X and with Y. You might expect $H(X,Y,Z)$ to be the three entropies minus the three mutual informations, which gives 3 bits. The true answer is 4, since Z adds no states. The formula counts twice what X tells about Y: X even tells us Z is even, which tells us Y is even, the same information X gives about Y directly. The correct formula, ``I believe,'' uses instead $I(X;Y|Z)$, what X tells about Y once Z is known, which here is zero. \nb{The formula holds for any three variables; it follows from the chain rule for entropy.}

I then show how to compute it: $I(X;Y|Z) = H(X|Z) + H(Y|Z) - H(X,Y|Z)$, where a conditional entropy is the entropy left after learning Z, averaged over Z's values. I mention two more theorems, $H(X|Y) = H(X,Y) - H(Y)$ and that zero $I(X;Z)$ and zero $I(Y;X|Z)$ give zero $I(X;Y)$, ``but I'm not going to go into those.'' \nb{Both theorems are correct.}

Mutual information is positive exactly when some $P(x,y) \neq P(x)P(y)$, that is, when $P(x|y) \neq P(x)$, which is what it means for one variable to be evidence about the other. Conditioned on Z, positive $I(X;Y|Z)$ means that even knowing Z, learning Y still changes our beliefs about X. Z screens off X from Y if, once Z is known, learning Y tells nothing more about X. For example, knowing Z is even, learning Y is 4 says nothing about whether X is 2, 4, 6 or 8. And ``without conditional independence, the universe would have no structure.'' I do not say why.

Now take five variables: speaks a language, has two arms and ten digits, wears clothes, can be poisoned by hemlock, has red blood. Each is evidence about the others. Learn that a thing, apple or rock or otherwise, speaks Chinese, and it more likely wears clothes. Learn that hemlock does not poison it, and it less likely has red blood. There are exceptions: Fred lost a finger in a volcano accident, Baby Barney does not talk yet, Irving the IRCBot talks but has no blood, and Nude Nellie talks but wears no clothes. So no one of the five screens off the rest. Knowing a thing is unclothed does not screen off what its speech says about its blood, since it may be Nude Nellie.

Five numbers that are all odd or all even would be simpler: knowing any one screens off the rest. Perhaps a sixth variable, ``human,'' would, even if we have to construct it rather than observe it. Pretending that it does is called ``Naive Bayes,'' and the name fits ``because it usually isn't quite true, but pretending that it's true can simplify the living daylights out of your calculations.'' Given ``human,'' Fred is no likelier to be a nudist, Nellie no less likely to speak, Barney no likelier to lack a limb. We do not track how clothing bears on speech given finger count. Instead, observations update the probability that the thing is human, or a chimpanzee or robot, and the class predicts the rest, such as hemlock. \nb{The pretence has a cost the post does not state. When observations tend to go together within the class, naive Bayes counts their shared evidence twice, as the three-variable formula did, and its probabilities can come out too extreme. Domingos and Pazzani (1997) found that its classifications can still be optimal when the assumption fails badly, while its probability estimates are optimal only when it holds.}

``Sound familiar?'' I show the blegg network from ``Neural Categories'': five property units, each linked only to a central category unit. With a logistic central unit and weights set by the likelihood ratios, it is ``exactly, mathematically equivalent to Naive Bayes.'' \nb{That is correct and standard.} It is ``a good guess'' that this is one reason logistic units work so well: they ``sneak in a little Bayesian reasoning while the designers aren't looking.'' \nb{No evidence is given. A trained logistic unit has the form of naive Bayes, but on data that break the independence assumption it need not end up with naive Bayes weights.}

So if someone presents a ``neural network'' covered in words like ``scruffy'' and ``emergent,'' do not assume it is ``Beyond the Realms of Logic.'' ``For this paradigm of adhockery, if it works, will turn out to have Bayesian structure.'' \nb{The only support in the post is its one example. In ``Searching for Bayes-Structure,'' posted two days before, the argument shows only that any useful part of a reasoning process must have ``at least a little Bayesian structure''; its stronger claim, that the search always finds the full structure, rests on the author's account of repeated discoveries.} The Bayesians will then explain exactly how the network works: its assumptions, the regularities it exploits, where it works and fails, and what its weights mean. ``Disappointing, isn't it?''
